\begin{agradecimentos}
\textit{Antes de mais nada, agradeço a Deus pela oportunidade de ingressar neste curso e nesta 
Universidade, onde encontrei algumas das melhores pessoas e amigos da minha vida. Mais do 
que isso, encontrei uma direção clara para o meu caminho. Hoje, sou feliz com o que faço e com 
a profissão que escolhi. Trabalhar com o que se ama não é um privilégio de todos, e sou 
imensamente grato por essa chance. Tenho um orgulho imenso de fazer parte desta instituição 
histórica e renomada, a Universidade de Brasília. Levarei comigo, para sempre, os momentos, 
os ensinamentos e as pessoas que aqui encontrei.}

\textit{Agradeço de forma especial à minha mãe, Dona Nininha, que sempre fez da minha educação 
sua prioridade, mesmo diante das dificuldades. Sou eternamente grato a essa mulher de coração 
manso e bondoso, a quem amo profundamente. Aos meus irmãos —  Neno e Teteus — que 
sempre estiveram ao meu lado, seja imprimindo material, me salvando com a minha moto 
quebradora (que, não raro, quebrava no caminho para a Universidade) ou simplesmente 
oferecendo uma palavra amiga.}

\textit{Gostaria de fazer um agradecimento especial à minha sobrinha Duda, que sempre foi uma fonte 
de inspiração para mim. O olhar dela, me vendo como alguém inspirador, sempre foi algo muito 
importante para mim. Esse sentimento me lembra o que sempre senti com o meu irmão Abadi, 
que também foi uma grande inspiração ao longo da minha vida e meu maior incentivador. 
Dentre os amigos que fiz aqui, não posso deixar de citar Maikão, Matheus (Bigode), Cauã, Caio 
Braga, Caio Lelis, Amanda, Malu, Pedro Jesus, JV, Matheus Soares e Rafael Nobre, que foram 
peças fundamentais nesta jornada. Agradeço de maneira especial ao Felipão, meu melhor 
amigo, com quem estou escrevendo este trabalho de conclusão de curso. Papiramos juntos no 
ensino fundamental, estudamos para o vestibular e para o concurso, e realizar este último 
trabalho do curso ao seu lado foi a concretização de um trato feito há muito tempo. Agradeço de 
coração por tudo, meu amigo.} 

\textit{E um agradecimento muito especial ao Ricardo, amigo do melhor tio do mundo, Clodoaldo 
(empatado tecnicamente com o tio Cláudio), que me apresentou o ATmega328P. Com aquele 
Arduino, fiz um robô se mover, com um código simples e tecnicamente horrível, mas que foi o 
suficiente para eu perceber, naquele momento, que eu já sabia o que queria fazer pelo resto da 
minha vida.}

\end{agradecimentos}
